\section{System Architecture and Design}
\label{sec:architecture}

The pipeline components described in this section are standard for late-fusion V2X systems. The one non-standard addition is self-beacon anchoring (Section~\ref{sec:beacon}), which enforces an ego-uniqueness invariant at the edge publish boundary. We describe the full pipeline here so that the failure modes characterized in Section~\ref{sec:eval} can be attributed to specific components.

To characterize the operational safety envelope of edge-assisted driving, we implement a representative late-fusion V2X pipeline and evaluate it as a closed-loop system. Our focus is the correctness of the state stream consumed by the planner under asynchronous delivery, partial ordering, and participant scaling. We instrument freshness-at-use (AoI at the moment the planner consumes edge state) and classify failures by their first terminating event: collision, ego-as-obstacle, other false-positive brake, or timeout.

Our system, shown in Fig.~\ref{fig:system_overview}, uses a late-fusion model. Late fusion keeps uplink bandwidth low and separates on-vehicle perception from edge-side state reconciliation, which allows us to study their interaction in isolation. We use a deterministic Kalman-filter-based tracker at the edge so that observed failures can be attributed to state-consistency effects rather than to non-determinism in learned fusion models. The following subsections detail the end-to-end architecture.

\subsection{Simulation Environment and World Model}
We build our system using the CARLA simulator~\cite{dosovitskiy_carla_nodate}, an open-source platform that provides high-fidelity models for vehicle physics, sensor data generation (including camera and LiDAR), and urban environments. All experiments run within CARLA, which is managed by our recently developed and open-sourced \ecav 
framework~\cite{eCAV}. Our framework allows for deterministic actor control and precise performance measurements.

In our scenarios, we differentiate between several types of road users. The \textbf{ego vehicle} refers to the primary Autonomous Vehicle (AV) whose decision-making and safety outcomes are being evaluated. Other participants can be \textbf{AVs} with the same perception and communication stack, or uninstrumented non-AV vehicles controlled by CARLA's traffic manager.

\subsection{Architectural Design: Late Fusion \& Hybrid Computation}
Our architecture uses a late-fusion approach. We opt for late-fusion (i.e., processing sensor data locally on each AV to generate object lists, then transmitting these lightweight lists) instead of early-fusion (transmitting raw sensor data), as it offers a scalable solution with significantly lower bandwidth demands, suitable for current and near-future V2X networks.


\subsection{On-Vehicle Perception and Tracking}

Each connected vehicle and Roadside Unit (RSU) acts as an intelligent sensor, responsible for processing its own raw sensor data to produce a lightweight, structured representation of its local view. All coordinates are transformed into a shared global coordinate frame before transmission to eliminate ambiguity at the fusion stage. This on-board stack includes a YOLOv5 object detector~\cite{glenn_jocher_2022_7347926}, a SORT tracker for 2D tracks~\cite{Bewley2016_sort}, and a LiDAR-based module for 3D localization. Local processing reduces the V2X data payload, a core principle of late-fusion architectures. A typical raw LiDAR scan can be megabytes in size, whereas a list of 3D bounding boxes is only a few kilobytes.

\subsection{V2X Data Transmission Protocol}

At each discrete time step of the simulation (50~ms), every agent transmits a structured data packet to the edge server. All data packets, both sensor detections and self-beacons, are tagged with a high-resolution generation timestamp from the sender. Our protocol includes two types of data from each vehicle:
\begin{enumerate}[label=\arabic*.,noitemsep,topsep=0pt,leftmargin=1.2\parindent]
    \item \textbf{Sensor-Detected Objects:} A list of 3D bounding boxes for all external objects detected by the on-board perception stack. Each box includes its dimensions (h, w, l), position (x, y, z), and yaw. 
    \item \textbf{Self-Beacon:} A single, authoritative record containing the vehicle's own precise ground-truth pose (location and orientation) and velocity at that time step, tagged with the vehicle's unique ID.
\end{enumerate}
The self-beacon provides the authoritative identity signal that the edge uses to enforce ego-uniqueness, as detailed in Section~\ref{sec:beacon}.

\subsection{Edge Intelligence Pipeline}
The edge pipeline reconstructs a consistent world model from multi-source, asynchronous inputs. Each packet carries a sender-generation timestamp on a common time base; the edge buffers packets and forms reconciled frames over a sliding window \([T,\,T{+}W)\), ordered by generation time rather than arrival time (Section~\ref{sec:timebase}). Within each reconciled frame, the edge must (i) suppress duplicate observations originating from overlapping fields of view and (ii) associate observations to persistent tracks despite delay variability and occasional reordering. While geometric deduplication and temporal alignment remove most duplicates, the tail of inter-source disagreement can still cause multiple tracks to persist for the same physical actor. When such duplicates involve the ego vehicle, the planner can misinterpret an ego duplicate as an external obstacle, triggering \emph{self-ghosting}. Section~\ref{sec:beacon} introduces self-beacon anchoring, which enforces an ego-uniqueness invariant at the edge and prevents ego duplicates from being published as obstacles.


\subsubsection{Data Aggregation and Deduplication}
The edge maintains a temporal buffer of incoming packets (detections and self-beacons). Before tracking, the edge applies lightweight deduplication that targets the dominant source of false multiplicity in late-fusion systems, namely overlapping detections of the same actor from different cameras and RSUs.

\paragraph{Perception-level cross-camera suppression.}
Each sensing agent performs cross-camera non-maximum suppression (NMS) to reduce duplicates before transmission. The suppression operates in the shared map frame and greedily removes near-identical boxes based on center proximity and box overlap. This step reduces redundant transmissions and prevents the tracker from seeing a large number of trivially duplicated boxes.

\paragraph{Tracker-level suppression and association gates.}
At the edge, AB3DMOT applies additional suppression and data association. Detections that remain after perception-level suppression are matched to existing tracks using 3D overlap-based association (Hungarian matching under a GIoU threshold) and track management rules (\texttt{min\_hits}, \texttt{max\_age}). This second-stage suppression is effective for the common case where inter-source disagreement is small.

\paragraph{Why deduplication is insufficient.}
Geometric suppression cannot eliminate duplicates by construction without sacrificing recall in dense scenes. When two nearby vehicles are physically close, aggressive suppression risks merging distinct actors, whereas conservative suppression admits a non-zero tail of duplicate ego detections due to depth uncertainty and cross-view disagreement. Our evaluation quantifies this tail and shows that it is sufficient to trigger self-ghosting under delay variability. Self-beacon anchoring addresses this residual by enforcing ego-uniqueness at publish time rather than relying solely on geometric heuristics.



\subsubsection{Fusion, Tracking, and Prediction}
After deduplication, the edge updates a persistent multi-object tracker and publishes a predictive world model.

\paragraph{Persistent tracker with jitter buffer.}
The edge runs AB3DMOT as a persistent tracker. Incoming packets are inserted into a jitter buffer indexed by generation time. At each dispatch time \(T\), the edge drains all packets whose generation timestamps fall in \([T,\,T{+}W)\) to form a reconciled frame. The edge processes self-beacons first (to update anchored ego state), then associates remaining detections to tracks under the standard AB3DMOT rules, and finally advances \(T\) monotonically. This design absorbs reordering and delay variability while avoiding costly repeated tracker re-initialization.

\paragraph{State and prediction outputs.}
For each stable track, the edge exports the current estimated state (pose and velocity) and a short-horizon prediction. Rather than fitting a least-squares line to the trajectory history, the predictor extrapolates using the Kalman filter's current velocity estimate, which converges to the true velocity within a few frames and avoids bias from noisy early observations at long range. The Kalman filter's measurement noise covariance is scaled up ($R \times 10$) to account for the camera-grade position error of $\pm 1$--$3$\,m typical of our pipeline, and the initial velocity covariance is inflated ($P_{7{:}} \times 10{,}000$) for aggressive early convergence. A track maturity gate suppresses collision predictions from immature tracks: anonymous tracks require at least 10 frames of history, while identified tracks require 3 frames. A displacement consistency check forces stationary predictions for tracks whose net displacement is small relative to their path length, preventing phantom velocity from position jitter. The edge broadcasts the resulting world model to participating vehicles.

Learned predictors such as Mamba~\cite{huang2025mambamotstatespacemodelmotion, mambatrack} can model non-linear motion but are harder to certify (ISO 26262)~\cite{Serna2020} and consume more of the edge compute budget~\cite{fusionlatencywacv2025}. We use a Kalman filter so that observed failures can be attributed to state-consistency effects rather than predictor non-determinism.

\subsubsection{Time base, synchronization, and reconciliation}
\label{sec:timebase}
Our design requires: (R1) per-message timestamp {generation}  from senders; (R2) a common time base across vehicles, RSUs, and edge; and (R3) event ordering via a sliding reconciliation window \(W\).

\textbf{Time source and distribution.}
On-vehicle clocks are disciplined by GNSS (PPS) and used by sensor drivers to stamp detections and self-beacons (Section \ref{sec:beacon}) at emission time. RSUs and edge servers run IEEE~1588 PTP on the backhaul; alignment errors are bounded and small relative to V2X jitter. The identifier carried by a self-beacon is orthogonal. Timestamps bind observations to a global time base for ordering.

\textbf{Event ordering.}
The edge executes an event-driven loop. At dispatch time \(T\), it forms a \emph{reconciled frame} containing all messages whose generation time lies in \([T,\,T{+}W)\), processes that frame (anchor/update beacons, then associate detections under the same gates and track rules), publishes the world model, and advances \(T\) monotonically. Packet reordering and variable network/compute delays are absorbed by \(W\).

\textbf{Holdover and degradation.}
If a participant loses GNSS/PTP lock, the edge admits its messages while the estimated drift remains below a small budget (holdover). Beyond that, cross-vehicle fusion for that participant is disabled and the ego falls back to planner-side admission until time sync is restored. The anchoring logic is unchanged.

\subsection{AV Control Loop and Asynchronous Data Fusion}
\label{sec:design:control-loop}
The ego vehicle's control loop operates on a periodic tick (every 50~ms in our simulation). In each tick, the onboard planner must make a decision based on the most current world model it can construct. This model is necessarily a fusion of two asynchronous data streams with different latencies:

\begin{enumerate}[label=\arabic*.,noitemsep,topsep=0pt,leftmargin=1.2\parindent]
    \item \textbf{Local Perception Data:} Information from the vehicle's own sensors. This data is fresh, with a latency determined by the onboard perception stack (e.g., ~25~ms in our tests).
    \item \textbf{Edge Intelligence Data:} The unified world model received from the edge. This data is inherently stale, with a total latency determined by the full vehicle-edge-vehicle round-trip (e.g., a realistic average of 190\,ms under the trace-driven model in \S\ref{sec:e2e_lat_model}).
\end{enumerate}

The AV planner does not halt its control loop to wait for edge data. In each tick, it fuses its latest local perception with the most recently received edge world model. The collision detection and prevention logic described below must tolerate the resulting temporal misalignment.
Safety enforcement operates at two layers.
At \textbf{fusion time}, the tracker enforces an Ego-Uniqueness invariant: at most one active track per persistent vehicle ID, with any delayed self-observation associated to the existing ego track rather than birthing a new obstacle track.
At \textbf{decision time}, the planner applies a staleness gate that discards edge contributions whose age-at-use exceeds a configured threshold. If all edge objects are rejected, the system degrades to local-only perception for that cycle.




\subsubsection{Collision Detection}

For each edge-predicted obstacle, the collision checker evaluates whether the obstacle's predicted trajectory intersects the ego vehicle's planned path within a configurable look-ahead horizon. Both trajectories are interpolated by arc length at 0.1\,m resolution and then reparameterized over time using each vehicle's speed. A collision is flagged when the spatial distance between the two time-indexed paths falls below a proximity threshold. The speed used for obstacle time reparameterization is derived from the predicted trajectory itself (displacement between the first two predicted points divided by the prediction time step), not from the detected trajectory, to reflect the predictor's estimate of future motion.

When the ego vehicle is stationary, its time-parameterized path is undefined, so a spatial-overlap fallback check is used instead. The ego vehicle's own track is excluded from collision candidates by matching on the vehicle's persistent ID, which is the only identity-based filter in the pipeline.

\subsubsection{Collision Prevention}

When a collision is predicted, the behavior agent enters an RSS-style proper response. On entry (predicted time-to-collision below the look-ahead horizon), the agent forces an emergency stop by asserting maximum braking ($\mathrm{brake} = 1.0$) rather than routing through the car-following controller, which would produce gentler deceleration insufficient for short stopping distances. The proper response remains latched until one of three formal exit conditions is met: (a) the ego vehicle has stopped (speed $< 2.0$\,km/h); (b) the threatening obstacle's lateral distance has been growing for three or more consecutive ticks, indicating cleared path; or (c) a time-to-live fallback expires. The latch prevents braking oscillation: without it, deceleration shifts the ego's time-reparameterized path such that the obstacle appears to pass safely, causing the planner to resume speed, which re-triggers the collision, producing repeated brake--accelerate cycles.






\section{Self-Beacon Anchoring for Latency Resilience}
\label{sec:beacon}



The edge publishes a world model that, by default, includes predictions for \textit{all} tracked objects, including the ego vehicle itself. Under delay and multi-source sensing, the ego's own track can appear as an external obstacle to its planner. We call the resulting false-braking failure \emph{self-ghosting}.

\textbf{Two-condition failure.}
Self-ghosting requires two conditions to hold simultaneously:
\begin{enumerate}[label=(\roman*),noitemsep,topsep=0pt,leftmargin=1.2\parindent]
\item \textbf{Ego-consistent duplicate:} the tracker maintains at least one track whose position and kinematics are consistent with the ego's state, in addition to or instead of the ego's own track.
\item \textbf{Missing identity label:} the ego-consistent track is unlabeled or mislabeled, so the planner cannot filter it as self.
\end{enumerate}
If either condition is absent, the planner does not treat the stale ego state as an obstacle. Multi-source stacks satisfy condition (i) when inter-camera position disagreement or delayed beacon arrival causes the tracker to birth a second ego-like track. Condition (ii) holds whenever the tracker lacks an authoritative identity signal to label the duplicate as ego.

\textbf{Mechanism and timeline.}
Self-ghosting arises when a late self-beacon births a stale track (T1), and a subsequent, fresher self-beacon cannot associate to T1 under the spatial gate, birthing a second track (T2) that the planner treats as a distinct obstacle. The sequence diagram in Fig.~\ref{fig:seq_selfghost} shows the causal timeline across the Ego, V2X network, and Edge tracker that can cause the emergence of the ``self-ghosting'' problem.

\input{floats/vector-self-ghosting}


\textbf{When does self-ghosting trigger?}
Let the edge predict from pose $P_0=(x,y,\psi)$ at $t_0$ (we omit $\psi$ below) with a constant-velocity model. A duplicate births when (i) the pose error at $t_1$ exceeds the gate, $\|\hat P_1-P_1\|>\tau$, and (ii) the stale track remains inside the ego's near-horizon reachable set so the planner treats it as actionable. Rapid \emph{decelerations} inflate $\|\hat P_1-P_1\|$ (overshoot), increasing duplicates; sustained \emph{accelerations} can reduce the error (undershoot) but do not eliminate the condition under arbitrary jitter. Anchoring removes this dependence by forcing ID-consistent association.

\textbf{Why simpler heuristics are insufficient.}
Proximity filters ($r$-gates), tracker hysteresis ($\tau,k$), source-tag seeding, and planner-only admission reduce, but do not eliminate duplicates. Each leaves residual false brakes in multi-actor scenes or sacrifices recall under occlusion. Anchoring yields a deterministic one-to-one identity mapping and removes the failure by construction.

\begin{algorithm}[t]
\footnotesize                      
\caption{Edge Data Fusion with Self-Beacon Anchoring}
\label{alg:edge_preprocess}
\begin{algorithmic}[1]
\Procedure{FuseDetections}{$\mathsf{AllDetections},\mathsf{FrameID}$}

  % ----------- STEP 1 --------------------------------------------------
  \Statex                                     % ← no line-number
  \Statex\textit{/* Step 1 — insert authoritative beacons */}
  \State $\mathsf{MasterList} \gets [\,]$
  \State $\mathsf{BeaconList} \gets \textsc{ExtractBeacons}(\mathsf{AllDetections})$
  \ForAll{$b \in \mathsf{BeaconList}$}
      \State $\mathsf{MasterList}\!\cdot\!\textsc{Add}(b)$
  \EndFor

  % ----------- STEP 2 --------------------------------------------------
  \Statex                                     % blank un-numbered spacer
  \Statex\textit{/* Step 2 — filter sensor detections */}
  \State $\mathsf{SensorList} \gets \textsc{ExtractSensors}(\mathsf{AllDetections})$
  \ForAll{$d \in \mathsf{SensorList}$}
      \If{\textsc{IsSliver}$(d)$}
          \State \textbf{continue}
      \EndIf
      \State $\textit{dup} \gets \textbf{false}$
      \ForAll{$b \in \mathsf{BeaconList}$}
          \If{$\textsc{IoU}(d,b) > \tau$}
              \State $\textit{dup} \gets \textbf{true};\ \textbf{break}$
          \EndIf
      \EndFor
      \If{\textbf{not}\ \textit{dup}}
          \State $\mathsf{MasterList}\!\cdot\!\textsc{Add}(d)$
      \EndIf
  \EndFor

  % ----------- STEP 3 --------------------------------------------------
  \Statex                                     % spacer again
  \Statex\textit{/* Step 3 — track */}
  \State \Return $\textsc{AB3DMOT.Track}(\mathsf{MasterList})$
\EndProcedure
\end{algorithmic}
\vspace{3mm}
\end{algorithm}

\textbf{Self-Beacon Anchoring:} We define a \emph{self-beacon} as a periodic, authenticated V2X message emitted by a vehicle that carries (i) the vehicle's persistent identifier and (ii) its current ego pose (position, orientation, velocity) in the shared map frame. Because the edge is the only component with a globally consistent view of all inputs at a given logical time step, we leverage this unique position to implement self-beacon anchoring. As shown in Algorithm~\ref{alg:edge_preprocess} and Fig.~\ref{fig:beacon_diagram}, the edge fusion logic is modified to \emph{prioritize} beacons. The ground-truth pose from a vehicle's self-beacon is used to create or update an authoritative ``anchor track'' for that vehicle within the AB3DMOT tracker.
All other sensor detections that spatially correspond to this anchor track are immediately fused into it, rather than creating a new, separate track. By anchoring each vehicle's identity to its own beacon at the point of fusion, the edge system tolerates delayed self-observations and avoids self-occlusion. Proximity filtering alone is insufficient in multi-vehicle scenes; the protocol enforces an Ego-Uniqueness Invariant at the system level. We modified the tracker to incorporate self-beacons as follows:

\begin{itemize}[noitemsep,topsep=0pt,leftmargin=1.1\parindent]

    \item \textbf{Persistent ID Propagation:} Third-party observers (RSUs, other AVs) detect anonymous objects without knowing their identity. The self-beacon carries both a persistent ID and a ground-truth pose. The edge uses spatial correlation to associate anonymous detections with the known beacon pose and propagates the persistent ID to the matched track.

    \item \textbf{Authoritative Matching:} We then modify the core data association logic. We introduce a high-priority rule that overrides the standard distance-based matching: if an incoming beacon and an existing track share the same persistent vehicle ID, they are forced to match, regardless of their spatial distance. This rule permanently associates a vehicle's beacon to its own track.

    \item \textbf{Specialized Track Management:} We alter the track life cycle logic. By default, trackers follow a ``birth-and-death'' model: a new track is born when a new object is detected, and it dies if it is not re-detected for a certain number of frames (max age). Our modification prevents an ego-vehicle's beacon from ever birthing a new, duplicate track. Furthermore, we make these authoritative beacon tracks more persistent by significantly increasing their max age, making them resilient to brief network dropouts.
    
\end{itemize}

\newcommand{\added}[1]{\textcolor{blue}{#1}}

\begin{algorithm}[t]
\footnotesize
\caption{AB3DMOT + Self‐Beacon Anchoring.}
\label{alg:beacon_logic}
\begin{algorithmic}[1]        % 1 = show line numbers

\Statex\textbf{Input:} Detection set $\mathsf{D}$, active track set $\mathsf{T}$

% ---------- STEP 0: split inputs ----------
\State $\mathsf{B}      \gets \textsc{ExtractBeacons}(\mathsf{D})$
\State $\mathsf{S}      \gets \textsc{ExtractSensors}(\mathsf{D})$
\State $\mathsf{M}      \gets [\,]$                           \Comment{matched pairs}
\State $\mathsf{U_T}    \gets \mathsf{T}$                     \Comment{unmatched tracks}

% ---------- STEP 1: beacon anchoring ----------
\Statex{/* Anchor beacons to tracks by persistent ID */}
\ForAll{$t \in \mathsf{T}$}                 % 1 indent level
    \ForAll{$b \in \mathsf{B}$}             % 2 indent levels
        \If{$t.\textit{id} = b.\textit{id}$}
            \State {$\mathsf{M}\!\leftarrow\!\mathsf{M}\cup(t,b)$}
            \State {$\mathsf{U_T}\!\leftarrow\!\mathsf{U_T}\setminus\{t\}$}
            \State {$\mathsf{B}\!\leftarrow\!\mathsf{B}\setminus\{b\}$}
            \State \textbf{break}
        \EndIf
    \EndFor
\EndFor

% ---------- STEP 2: standard association ----------
\State $(\mathsf{M}',\mathsf{U_T},\mathsf{U_D}) \gets
       \textsc{Associate}(\mathsf{S},\mathsf{U_T})$
\State $\mathsf{M}\!\leftarrow\!\mathsf{M}\cup\mathsf{M}'$

% ---------- STEP 3: update matched ----------
\ForAll{$(t,d)\in\mathsf{M}$}
    \State $t.\textsc{Update}(d)$
\EndFor

% ---------- STEP 4: create tracks for unmatched dets ----------
\ForAll{$d \in \mathsf{U_D}$}
    \If{\textbf{not}\,d.\textit{isBeacon}}
        \State \textsc{CreateTrack}($d$)   \Comment{ghost-prevention}
    \EndIf
\EndFor

% ---------- STEP 5: age & prune ----------
\ForAll{$t \in \mathsf{U_T}$}
    \State $t.\textit{age}\!\gets\!t.\textit{age}+1$
    \If{$t.\textit{isBeacon}$ \textbf{and} $t.\textit{age} > \text{MAX\_AGE}_B$}
        \State \textsc{DeleteTrack}($t$)
    \ElsIf{\textbf{not}\,t.\textit{isBeacon} \textbf{and} $t.\textit{age} > \text{MAX\_AGE}$}
        \State \textsc{DeleteTrack}($t$)
    \EndIf
\EndFor

\State \Return $\mathsf{ActiveTracks}$
\end{algorithmic}
\end{algorithm}

\begin{figure}[t]
    \centering
    \includegraphics[
        trim={0cm 5.3cm 3.6cm 5.3cm}, 
        clip,
        width=\columnwidth % Example: Makes the image 90% of the full page width
    ]{figures/self_beaconing.pdf} 
    \caption{The Self-Beacon Anchoring mechanism.}
    \label{fig:beacon_diagram}
\end{figure}

Our beacon logic modifications are shown in Algorithm \ref{alg:beacon_logic}. 
\paragraph{Identity provenance and Ego-Uniqueness.}
Self-beacons carry a verifiable vehicle identifier over a bounded anchoring epoch. The edge maintains at most one active track per identifier and enforces an ID-consistent association (forced match): on a beacon with ID \(A\), update the anchored track for \(A\) even if geometric gating would reject.

\paragraph{Temporary ID rotation.}
To decouple the anchoring protocol from long-lived identifiers, the edge assigns each vehicle a random temporary ID at session start through a Beacon ID Manager. The tracker operates on temporary IDs internally and reverse-maps to the vehicle's session identity at publish time. Temporary IDs can be rotated periodically (analogous to BSM pseudonym changes under SAE~J2945) without disrupting track continuity, since the Beacon ID Manager maintains the mapping across rotations. The anchoring invariant depends only on the presence of a valid identifier at use time; certificate issuance and revocation follow standard V2X PKI procedures and are orthogonal to the protocol logic.





